\begin{multicols}{2}

\chapter{Science Fiction Bestiary}\index{Creatures and NPC!Science Fiction}

\section{Basic Creatures and NPC for a Science Fiction Game}\index{Creatures and NPC!Science Fiction!Basic}
\index[AllCre]{Short listings!Science Fiction}

\textbf{Innocuous rodent}: level 1

\textbf{Guard beast}: level 3, perception as level 4

\textbf{Corporate drone}: level 2

\textbf{Physical laborer}: level 2; health 8

\section{Science Fiction Creatures and NPCs by Level and Tech}\index{Creatures and NPC!Science Fiction!By level}

\end{multicols}

\newpage

\begin{table}

\centering

\caption{Science Fiction Creatures and NPCs by Level and Tech}
\label{tab:Science Fiction Creatures and NPCs by Level and Tech}

\begin{tabularx}{\textwidth}{| Y | Y | Y |}
\hline
\textbf{Level} & \textbf{Name} & \textbf{Tech Rating} \\ \hline

1 & Space rat & Advanced \\ \hline
2 & Silicon parasite & Advanced \\ \hline
3 & Infovore & Fantastic \\ \hline
3 & Mock organism & Advanced \\ \hline
3 & Natathim (homo aquus) & Advanced \\ \hline
3 & Sentinel tree & Advanced \\ \hline
3 & Zero-point phantom & Fantastic \\ \hline
4 & Devolved & Advanced \\ \hline
4 & Ecophagic swarm & Advanced \\ \hline
4 & Hungry haze & Fantastic \\ \hline
4 & Inquisitor & Fantastic \\ \hline
4 & Malware, fatal & Advanced \\ \hline
4 & Redivus & Fantastic \\ \hline
4 & Wraith (homo vacuus) & Advanced \\ \hline
5 & SHining one & Fantastic \\ \hline
5 & Synthetic person & Advanced \\ \hline
6 & Vacuum fungus & Advanced \\ \hline
6 & Exoslime & Fantastic \\ \hline
6 & Photonomorph & Fantastic \\ \hline
7 & Posthuman 7 Fantastic \\ \hline
7 & Thundering behemoth & Fantastic \\ \hline
8 & AI & Advanced \\ \hline
8 & Cybrid & Fantastic \\ \hline
8 & Wharn interceptor & Fantastic \\ \hline
10 & Godmind & Fantastic \\ \hline
10 & Omworwar & Fantastic \\ \hline

\end{tabularx}
\end{table}

\begin{multicols}{2}

\section{Science Fiction Bestiary}\index[SFCre]{Science Fiction Bestiary}
\index{Creatures and NPC!Science Fiction!Bestiary}

\subsection{Artificial Intelligence (AI) 8 (24)}\index[SFCre]{Science Fiction Bestiary by name!Artificial Intelligence}
\index[SFCre]{Science Fiction Bestiary by level!8!Artificial Intelligence}
\index[AllCre]{Creatures by name!Artificial Intelligence}\index[AllCre]{Creatures by level!8!Artificial Intelligence}

If a supercomputer can think independently, it’s a strong AI (an artificial intelligence). Though not as advanced as godminds, AIs can develop inscrutable goals. 
AIs take many forms. Some are distributed across a vast network. Others are encoded into a singular “computer core.” A few are machines with organic parts. All are entities of extreme intelligence able to adapt to new situations, and most act on some kind of plan, whether long-acting, or newly concocted to fit the situation at hand. 

\textbf{Motive}: Varies 

\textbf{Environment}: Almost anywhere 

\textbf{Health}: 33 

\textbf{Damage Inflicted}:  10 

\textbf{Armor}: 2 

\textbf{Movement}: Immediate 

\textbf{Modifications}: Speed defense as level 2, knowledge tasks as level 9 

\textbf{Combat}: An electrical discharge—or in some cases precisely pulsed sequences of lights, each designed for a specific creature to see—can affect all targets within short range of the AI (or the AI’s local terminal), inflicting 10 points of damage from electricity (or 10 points of Intellect damage, which ignores Armor). 
Some AIs can take an action to absorb matter around them (such as walls, floor, equipment, unresisting living creatures, and so on), regaining 5 points of health. 
An AI is likely able to deploy cyphers and artifacts in combat and also relies on guardians (such as synthetic people made to its own design) to aid it. Unless a particular AI uses a computer core, damage to an AI may just be damage done to a “terminal,” so even if an AI is seemingly destroyed, it might exist as another instance somewhere else. 

\textbf{Interaction}: Some AIs enjoy negotiation. Others simply ignore humans as unworthy of their time and attention. An AI’s voice often sounds surprisingly human. 

\textbf{Use}: The characters are contacted by an AI sympathetic to biological beings. It wants them to accomplish a task on a moon of Jupiter: assassinate a security officer who the AI calculates as being a nexus of future disaster if he isn’t removed from the equation. 

\textbf{Loot}: An AI might have access to 1d6 cyphers and possibly an artifact or two.

\textbf{GM Intrusion}: The AI knows a phrase and series of images to flash at a particular PC to stun them for around as it attempts to upload an instance of itself into their mind.

\subsection{Cybrid 8 (24)}\index[SFCre]{Science Fiction Bestiary by name!Cybrid}
\index[SFCre]{Science Fiction Bestiary by level!8!Cybrid}
\index[AllCre]{Creatures by name!Cybrid}\index[AllCre]{Creatures by level!8!Cybrid}

Cybrid origins could be the result of someone finding a cache of ancient ultra technology, or manufactured by a post-singularity AI for some unfathomable purpose, or even the result of banned weapons research by a nation-state or conglomerate. The human remnants in each cybrid’s carbon fiber and nested shells of nanotech exist in a red haze of pain; neuro-wetware and chemicals bathing their remaining living tissues hold the pain partly at bay. 

From the exterior, not much of the original human is obvious, except perhaps in the echo of a humanoid shape. Each one has a unique conformation, but all are designed to strike fear in anyone seeing one, ally and enemy alike. 

\textbf{Motive}: Kill away the pain 

\textbf{Environment}: Usually set to guard important areas, creatures, or objects, or deployed in war 

\textbf{Health}: 60 

\textbf{Damage Inflicted}:  10 

\textbf{Armor}: 3 

\textbf{Movement}: Short; flies a very long distance each round; can maneuver like an autonomous level 5 spacecraft if using extended vehicular combat rules. 

\textbf{Combat}: Cybrids can attack up to three foes that they can see up to about 300 m (1,000 feet) away as a single action with graser (gamma ray laser) beams, inflicting 10 points of damage on each target and everything in immediate range of the target. Those caught in the beam who succeed on a Speed defense roll still suffer 2 points of damage. If the cybrid focuses on a single target, treat the attack as a level 10 attack that inflicts 14 points of damage, or 6 points even on a successful Speed defense roll.

Self-repair mechanisms allow the creature to regain 2 points of health per round. 

\textbf{Interaction}: If communication can be opened up through a cybrid’s haze of pain, it might be possible to temporarily wake the consciousness of the human remnant inside. However, that remnant consciousness might not be happy to discover what it’s become. 

\textbf{Use}: A cybrid has appeared in orbit around the station, ship, or moon with a compromised life support system or fragile dome. If it engages, the death toll will be staggering. 

\textbf{Loot}: PCs who investigate the inert remains of the creature discover several manifest cyphers

\textbf{GM Intrusion}: The character struck by the graser beam develops radiation poisoning, in this case a level 8 disease that drops the character one step on the damage track each day that it goes untreated..

\subsection{Devolved 4 (12)}\index[SFCre]{Science Fiction Bestiary by name!Devolved}
\index[SFCre]{Science Fiction Bestiary by level!4!Devolved}
\index[AllCre]{Creatures by name!Devolved}\index[AllCre]{Creatures by level!4!Devolved}

Conglomerate security subsidiaries regularly experiment with new ways to create super-soldiers, either to supply to a government on a contract basis, or to use for themselves. These experiments produced hundreds of dead ends—literally—plus a few dangerous failures. The devolved are one of those dangerous failures. These malformed, hideous brutes share a common heritage but display a wide array of maladies and mutations in the flesh, including withered limbs or elephantine patches of thick, scaly skin, misplaced body parts, and mental abnormalities. Simple-minded and afflicted with pain from their twisted, broken forms, the devolved vent all their hatred and wrath against all others. 

Even successfully created super-soldiers require a regular regimen of specialized drugs to keep them healthy. Most are shipped out to fight on faraway fronts, whether that’s on a distant space station, moon, or in another star system entirely. Without their drugs, they may devolve.

\textbf{Motive}: Hungers for flesh 

\textbf{Environment}: Groups of three to five, usually in locations where organized security can’t easily reach 

\textbf{Health}: 21 \textbf{Damage Inflicted}: 6 to 12 points 

\textbf{Movement}: Short 

\textbf{Modifications}: Intimidation tasks as level 6; Intellect defense and Speed defense as level 2 due to malformed nature 

\textbf{Combat}: Devolved attack with a claw, a bite, or some other body part, inflicting 6 points of damage. They throw themselves at their enemies with mindless ferocity and little regard for their own safety. Easily frustrated, a devolved grows stronger as its fury builds. Each time it misses with an attack, the next attack is eased by one additional step and the damage it inflicts increases by 2 points (to a maximum of 12 points). Once the devolved successfully inflicts damage on a target, the amount of damage it inflicts and the difficulty of its attacks returns to normal. Then the cycle starts anew. 

\textbf{Interaction}: Devolved speak when they must, punctuating their statements with growls and barks. Their understanding seems limited to what they can immediately perceive, and they have a difficult time with abstract concepts. 

\textbf{Use}: An expedition to a ruined conglomerate research facility uncovers a cyst of devolved that live within its sheltering bunkers. 

\textbf{Loot}: For every three or so devolved, one is likely to carry a cypher

\textbf{GM Intrusion}: The devolved detonates upon its death, inflicting 6 points of damage on everything in immediate range.

\subsection{Ecophagic Swarm 4 (12)}\index[SFCre]{Science Fiction Bestiary by name!Ecophagic Swarm}
\index[SFCre]{Science Fiction Bestiary by level!4!Ecophagic Swarm}
\index[AllCre]{Creatures by name!Ecophagic Swarm}\index[AllCre]{Creatures by level!4!Ecophagic Swarm}

Tiny nanomachines can be incredibly useful tools. But they can also become a terrible threat. Like cells in a living body that develop cancer, these out-of-control self-replicating robots can consume everything in their path while building more of themselves. A typical swarm is about 6 m (20 feet) in diameter, individually consisting of millions of individual minuscule machines. However, several swarms can act together, creating a much larger cloud of death with just one purpose: to eat and replicate. Able to move large distances by gliding through the air, cloud-like swarms take on intriguing shapes and ripple with mathematical patterns as they approach a potential target, beautiful and deadly. 

Ecophagic swarms sometimes build weird structures or artifacts in the wake of their feeding, like massive metallic ant or wasp mounds, or something without any reference at all in the natural world.

\textbf{Motive}: Hungers for matter, including flesh 

\textbf{Environment}: Ecophagic swarms are drawn most to areas rich in rare-earth metals, such as large cities or space stations where everyone carries a smartphone, AR glasses, or something similar 

\textbf{Health}: 12 

\textbf{Damage Inflicted}:  4 points 

\textbf{Movement}: Flies a long distance 

\textbf{Combat}: As a mass of countless tiny machines, an ecophagic swarm can flow around obstacles and squeeze through cracks large enough to permit a single sub-millimeter machine. That includes over and around other creatures. Characters touched by a leading edge—or wholly enveloped within the hazy “body”—of an ecophagic swarm must succeed on a Might defense task or take 4 points of damage. If the character doesn’t wear armor of some kind, they take 1 point of damage even if they succeed. 

For its part, an ecophagic swarm ignores any attack that targets a single creature (unless it’s an electrical attack), but it takes normal damage from attacks that affect an area (and electrical attacks), such as a detonation. A swarm cannot enter liquids, unless it takes about an hour to build new subunits that are aquatic. 
\textbf{Interaction}: Someone with an ability to communicate with machines might be able to interact with a swarm. Even then, attempts to influence it are hindered by three steps. 

\textbf{Use}: A promising new nanotech “printing” technology was hacked by radical elements

\textbf{GM Intrusion}: The character must succeed on a Speed defense roll or their armor (or other important piece of equipment) is taken by the swarm.

\subsection{Exoslime 6 (18)}\index[SFCre]{Science Fiction Bestiary by name!Ecoslime}
\index[SFCre]{Science Fiction Bestiary by level!6!Ecoslime}
\index[AllCre]{Creatures by name!Exoslime}\index[AllCre]{Creatures by level!6!Exoslime}

Amoeboid life predominates in some environments. Sometimes, it slimes asteroid crevices or its greasy residue is found on abandoned spacecraft. In a few cases, large portions of entire worlds are covered in living seas of translucent protoplasm. Individual volumes of exoslime are 5 m (15 foot) diameter moldlike blobs. Exoslimes possess independent minds, but in some settings may be manufactured entities designed to explore new locations, interact with aliens, or subjugate aliens. Exoslimes can learn to respect the autonomy of other creatures, though their natural instinct is to absorb novel objects and creatures they discover in order to learn about them. Exoslimes can also replicate anything they absorb, even a previously eaten living intelligent being. 

\textbf{Motive}: Hungers for information 

\textbf{Environment}: Moist and warm areas 

\textbf{Health}: 33 

\textbf{Damage Inflicted}:  6 points 

\textbf{Movement}: Immediate; immediate when climbing or burrowing 

\textbf{Modifications}: Speed defense as level 5 due to size 

\textbf{Combat}: Though slow, an exoslime is dangerous. When roused, all characters within immediate range of an exoslime must succeed on a Might defense roll each round or be touched by the heaving mass. A victim adheres to the slime’s surface and takes 6 points of acid damage each round. The victim must succeed on a Might defense roll to pull free. A victim who dies from this damage is consumed by the exoslime. The exoslime may later create a duplicate of any previously devoured fleshy creature, a process requiring about three rounds to complete. Duplicates have full autonomy, and can communicate with the slime. 

\textbf{Interaction}: An exoslime prefers to eat a newly-encountered creature, then create a duplicate of it to act as a translator. Of course, a stranger might not understand why the exoslime is trying to eat it. 

\textbf{Use}: The sample brought in from the exterior has a weird, mucus-like growth that seems able to slowly eat through most materials.

\textbf{GM Intrusion}: The character escapes an exoslime attack, but a piece of quivering protoplasm remains stuck to their flesh, eating away at 1 point of Speed damage (ignores Armor) each round until the character succeeds on a Might roll as an action.

\subsection{Godmind 10 (30)}\index[SFCre]{Science Fiction Bestiary by name!Godmind}
\index[SFCre]{Science Fiction Bestiary by level!10!Godmind}
\index[AllCre]{Creatures by name!Godmind}\index[AllCre]{Creatures by level!10!Godmind}

Unfathomably powerful post-singularity AIs, godminds are vast, having used the matter of an entire solar system and all its planets to create an immense brain, weave themselves into a nebula, or encode themselves into quantum strings of existence light-years across. When necessary, a godmind forms a nexus of consciousness—an instance—appearing as a disembodied eye of electromagnetic energy, ranging from about the size of a human eye all the way up to the size of a planet. 

\textbf{Motive}: Ineffable 

\textbf{Environment}: Anywhere, usually in space 

\textbf{Health}: 50 (per instance) 

\textbf{Damage Inflicted}:  15 points 

\textbf{Movement}: Very long when flying 

\textbf{Combat}: A godmind can vary the physical laws of the universe within a light-second of one of its instances (some would call them avatars) to create an effect most useful to the godmind at the time. For instance, a godmind could create a gamma ray burst inflicting 15 points of damage on all creatures within very long range, attempt to put a target into temporal stasis, send a target (even a target as large as spacecraft) through a temporary wormhole gate, and so on. It could also scan the memory banks of any digital machine, and possibly of any living creatures. In any event, if an instance were targeted, and successfully neutralized or even destroyed, the godmind itself isn’t harmed. An aggressor would have to find the godmind’s primeval “computer core” to destroy one, likely an epic quest in and of itself. 

\textbf{Interaction}: To actually get a godmind’s attention and negotiate could require ancient command code, finding an old input device, or showing up with a relic from an ancient ultra or other prize. If a godmind does render aid, it’s likely to be in a form that is initially enigmatic, though ultimately extremely powerful. 

\textbf{Use}: A universal threat requires a defense that is equally potent. Research suggests that the diffuse nebula known as the Double Helix may actually be the visible form of a vast godmind. Perhaps it can help.

\textbf{Loot}: Sometimes a godmind provides powerful artifacts to aid those who petition them for aid, assuming the need is dire.

\textbf{GM Intrusion}: The godmind rewinds time a few seconds and sidesteps whatever negative effect would have otherwise
inconvenienced it.

\subsection{Hungry Haze 4 (12)}\index[SFCre]{Science Fiction Bestiary by name!Hungry Haze}
\index[SFCre]{Science Fiction Bestiary by level!4!Hungry Haze}
\index[AllCre]{Creatures by name!Hungry Haze}\index[AllCre]{Creatures by level!4!Hungry Haze}

Hungry hazes are found in regions where the fundamental laws of physics have been eroded or are weak. They are named for how they appear as distortions of sight, like areas of heat haze, that shimmer in the air. These colorless hazes rapidly advance when they sense prey, taking on a “hungry” orange-red hue as they cling to the bodies of whatever they attempt to feed on next. 

Victims being fed upon by a hungry haze sometimes hallucinate, seeing a physically manifest monster instead of formless vapor.

\textbf{Motive}: Hungers for flesh 

\textbf{Environment}: Alone or in groups of three to five, usually in areas of strained space-time. Immune to the effects of vacuum. 

\textbf{Health}: 12 

\textbf{Damage Inflicted}:  5 points 

\textbf{Movement}: Flies an immediate distance each round 

\textbf{Modifications}: Stealth tasks as level 5 

\textbf{Combat}: A hungry haze breaks down the flesh of all living creatures within immediate range, inflicting 5 points of damage. As an insubstantial haze, only attacks that affect an area have a chance to inflict full damage on them; other successful attacks only inflict 1 point of damage, regardless of the amount indicated. If a hungry haze successfully feeds, it gains 1 point of health, even if the increase puts it above its maximum health. If a hungry haze is reduced to zero health, a smooth thumb-sized egg of unknown material is left behind. 

\textbf{Interaction}: A hungry haze does not speak or seem to have language. But it is not mindless; it can learn from its experiences and figure out creative solutions to problems. 

\textbf{Use}: After a research station on Mercury is abandoned for unspecified issues, salvagers show up looking for easy pickings. But a strange haze seems to hang over the station. 

\textbf{Loot}: People (or AI) interested in strange manifestations would probably pay for the remains of a hungry haze in an amount equal to the expensive price category.

\textbf{GM Intrusion}: The character’s Armor rating is reduced by 1; the hungry haze apparently can eat more than just flesh.


\subsection{Infovore 3 (9)}\index[SFCre]{Science Fiction Bestiary by name!Infovore}
\index[SFCre]{Science Fiction Bestiary by level!3!Infovore}
\index[AllCre]{Creatures by name!Infovore}\index[AllCre]{Creatures by level!3!Infovore}

Entities of information with an affinity for technology, infovores are nothing but stored information without a bit of mechanism to inhabit. But once one gains control of a device, computer system, or other powered item, it self-assembles over the course of a few rounds, becoming stronger and more dangerous as each second passes. Luckily, an infovore seems unable to hold this form for long, and whether defeated or not, it eventually falls back into so much scattered junk. But in one of those objects, the core of the infovore remains, waiting to come into close enough proximity to another fresh mechanism to begin the rebirth process again. 

Infovores have also been called ghost fabricators and aterics

\textbf{Motive}: Hungers for information 

\textbf{Environment}: Anywhere powered devices are found 

\textbf{Health}: 9 

\textbf{Damage Inflicted}:  3–10 points 

\textbf{Armor}: 3 

\textbf{Movement}: Short 

\textbf{Modifications}: Attacks and defends at an ever-escalating level 

\textbf{Combat}: A newly animate infovore (level 3) has a rough but articulated form that it uses to batter and cut targets who carry powered devices on them. Unless destroyed, on each subsequent round it draws nearby inert mechanisms, unattended metallic and synthetic matter, and ambient energy, and its effective level increases by one. This level advancement completely heals all previous damage it has taken and advances it to the amount of health consistent with a creature of the next higher level. Damage, attacks, and defense continue to ramp up as well, continuing each round until the creature is either destroyed or it reaches level 10. After being active for one round at level 10, it spontaneously disassembles, falling back into so many scattered pieces of junk. Finding the “seed” device amid this junk is a difficulty 6 Intellect-based task. 

\textbf{Interaction}: Infovores are fractured, fragmented beings. Characters who can talk to machines might be able to keep one from “spinning up” to become a threat and learn something valuable, but only for a short period. 

\textbf{Use}: Among the devices collected from trade, salvage, archeological dig, or some other unique source, one was actually an inactive infovore, quiescent until plugged in or scanned. 

\textbf{Loot}: An infovore that has undergone spontaneous disassembly leaves one or two manifest cyphers; however, there’s a chance that one of those cyphers is actually the infovore seed.

\textbf{GM Intrusion}: The character must succeed on a Speed defense task or lose a powered piece of equipment (an artifact) or a manifest cypher as it’s pulledinto the self-assembling infovore. The infovore gains an additional attack each round.

\subsection{Inquisitor 4 (12)}\index[SFCre]{Science Fiction Bestiary by name!Inquisitor}
\index[SFCre]{Science Fiction Bestiary by level!4!Inquisitor}
\index[AllCre]{Creatures by name!Inquisitor}\index[AllCre]{Creatures by level!4!Inquisitor}

Inquisitors are aliens who call themselves “inquisitors” when they contact new species. Their preferred method of interaction is to study a given area for its flora and fauna, and attempt to collect a representative sample of any intelligent species they find (such as humans). Collected subjects may be gone for good, but other times they wake with little or no recollection of the experience save for bruises, missing digits or teeth, scabbed-over circular head wounds, and a gap of three or more days in their memory. Instead of arms, inquisitors sprout three sets of three tentacles like those of a squid, each of which branches into a smaller and finer set of manipulator tendrils. They can manipulate complex machines in a way that a regular human could never hope to. In most settings, inquisitors possess a level of technology and advancement well above that enjoyed by humans. 

\textbf{Motive}: Knowledge 

\textbf{Environment}: In groups of three to twelve 

\textbf{Health}: 18 

\textbf{Damage Inflicted}:  6 points 

\textbf{Movement}: Short; short when climbing 

\textbf{Modifications}: Knowledge-related tasks as level 8 

\textbf{Combat}: Inquisitors can batter and squeeze foes with their tentacles, but they prefer to use advanced items that they always carry, including long-range energy weapons that can inflict damage or, with a flipped setting, induce deep sleep for an hour or more if the victim fails a Might defense task. Usually, inquisitors attempt to cause as little damage as possible to potential subjects, so the sleep setting is used most often. They also carry defensive items, including manifest cyphers that can grant +4 to Armor for a few minutes or throw up a level 8 force field barrier. In case a specimen collection mission goes badly, at least one inquisitor carries a manifest cypher that creates a short-lived teleportation portal for instant transport to a distant and hidden base (which might be a spacecraft or a transdimensional redoubt). 

\textbf{Interaction}: Inquisitors are always eager to “talk,” though they usually end up wanting to know a lot more than characters are willing to divulge. 

\textbf{Use}: An entire freehold on Mars goes missing. Left-behind clues point to inquisitors. 

\textbf{Loot}: Most inquisitors carry a couple of manifest cyphers that have offensive and defensive capabilities.

\textbf{GM Intrusion}: The character (or characters) wake after a long rest, only to realize that more than ten hours have passed. They all have strange marks and wounds, but no one remembers why. One character—an NPC or follower—might even be missing.

\subsection{Malware, Fatal 4 (12)}\index[SFCre]{Science Fiction Bestiary by name!Malvare, Fatal}
\index[SFCre]{Science Fiction Bestiary by level!4!Malvare, Fatal}
\index[AllCre]{Creatures by name!Malvare, Fatal}\index[AllCre]{Creatures by level!4!Malware, Fatal}

This purely malefic program has aggressive machine learning capabilities, allowing it to accomplish truly innovative and nasty tricks. Fatal malware may have originated as a simple virus or spyware coded for a specific purpose, but corruption and lightning-quick electronic evolution has turned it into something that exists purely to infect orderly electronic systems, spacecraft, space stations, smart weapons, and anything else with an operating system. Infected objects turn against living people. An instance often has the form of the system it’s infected, but occasionally fatal malware physically manifests as a metallic “cancer” of wires and self-assembling circuits hanging like a tumor across a server room, shipmind core, or data center, having perverted the original machine’s self-repair functions. Sometimes 4D printers are also compromised. 

\textbf{Motive}: Corruption and destruction 

\textbf{Environment}: Any electronic system able to run code can host one or more instances 

\textbf{Health}: 18 

\textbf{Damage Inflicted}:  5 points 

\textbf{Movement}: As the system it infects 

\textbf{Modifications}: Knowledge tasks related to computers and other electronic systems as level 6 

\textbf{Combat}: An instance of fatal malware that physically touches (or electrically connects with) a powered device of up to level 6 can attempt to seize control of it. It can then use that device to attack living targets. If the controlled system is a computer, smartphone, AR glasses, or some other piece of equipment that doesn’t have any intrinsic movement, the malware attempts to electrocute a user, or if a smart weapon, cause some kind of fatal accident with it. A compromised computer or shipmind voice can dangerously mislead victims. Fatal malware duplicates itself, creating many instances, and those that survive are usually slightly better at avoiding being erased than the previous generations. 

\textbf{Interaction}: Fatal malware isn’t really sentient and thus can’t really be negotiated with; some instances could mimic intelligence to draw humans into a trap. 

\textbf{Use}: An instance of fatal malware has gotten into a shipmind, which is making the normally trustworthy AI act out in unexpectedly dangerous ways. The shipmind itself doesn’t know it’s infected.

\textbf{GM Intrusion}: The fatal malware divides into a second instance and attempts to override and control another piece of equipment carried by the character, especially a character with cybernetic implants.

\subsection{Mock Organism 3 (9)}\index[SFCre]{Science Fiction Bestiary by name!Mock Organism}
\index[SFCre]{Science Fiction Bestiary by level!3!Mock Organism}
\index[AllCre]{Creatures by name!Mock Organism}\index[AllCre]{Creatures by level!3!Mock Organism}

Artificial life can be created by selective breeding, synthetic and genetic engineering, or by accidental miscalculation in some unrelated high-energy or food-research program. When artificial life takes a wrong turn, the results run the gamut from disappointing to dangerous. If an artificial entity starts out benign, it’s difficult to know if a hidden or slowly developing flaw will tip it over the edge into dangerous dysfunction—or if it just acts oddly because it doesn’t know the social cues. Should synthetic beings be treated as people, pets, or monsters to be stamped out and destroyed? That’s the eternal question and one that’s usually answered by those most afraid of potential dangers that might accompany the creation of something no one intended. 

\textbf{Motive}: Defense or destruction 

\textbf{Environment}: Usually in secluded locations alone unless hiding in unused storage rooms of a large facility 

\textbf{Health}: 18 

\textbf{Damage Inflicted}:  5 points 

\textbf{Armor}: 2 

\textbf{Movement}: Short 

\textbf{Combat}: A mock organism can release an electrical discharge against a target at short range. In melee, a mock organism’s poisoned claws inflict damage and require the target to succeed on a Might defense task, or the poison induces a coma-like slumber in the target. Each round the target fails to rouse—an Intellect task—they take 3 points of ambient damage. 

\textbf{Interaction}: A mock organism is intelligent and can sometimes be swayed by reason. It might be passive, but if disturbed in a place it thought was secure against intrusion, it could grow belligerent and even murderous. Once so roused, a mock organism might still be calmed, but all such attempts are hindered. 

\textbf{Use}: A scientist’s ruined lab contains several unexpected surprises, including a mock organism that yet grieves over the loss of its creator. 

\textbf{Loot}: A mock organism requires many parts. Salvage from a destroyed mock organism could result in a manifest cypher or two and another item that, with a bit of jury-rigging, works as an artifact. 

\textbf{GM Intrusion}: The character hit by the mock organism’s melee attack doesn’t take normal damage. Instead, the mock organism drops onto the character. The PC is pinned until they can succeed on a difficulty 6 Might-based task to escape. While pinned, the creation whispers mad utterances into the target’s ear.

\subsection{Natathim 3 (9)}\index[SFCre]{Science Fiction Bestiary by name!Natathim}
\index[SFCre]{Science Fiction Bestiary by level!3!Natathim}
\index[AllCre]{Creatures by name!Natathim}\index[AllCre]{Creatures by level!3!Natathim}

Genetically engineered to live in the water oceans discovered beneath the ice crusts of various solar moons, natathim (Homo aquus) have human ancestors, but barely look it. Survival in the frigid, lightless depths of extraterrestrial oceans required extreme adaptation. Predominantly dark blue, their undersides countershade to pure white. Though humanoid, their physiology is streamlined, giving their heads a somewhat fish-like shape, complete with gills and large eyes to collect light in the depths. Their bodies are adorned with fins and frills, including a long shark-like tail, and they have webbed extremities with retractable claws. 

Depending on the setting, natathim are either human allies with the same (or even more advanced) tech, enemies with the same or more advanced tech, or genetic anomalies treated like laboratory rats burning with genocidal fury at what’s been done to them. Alternatively, natathim could be discovered in Earth’s deepest oceans, their origin mysterious, but able to interbreed with humans as a method for maintaining their line. 

\textbf{Motive}: Just as with humans, natathim have many and varied motivations and drives. 

\textbf{Environment}: Anywhere in or near water, or in suits/craft with marine environments, in schools of three to twelve. Natathim can act normally in air for up to twenty-four hours before they must return to water. 

\textbf{Health}: 9 

\textbf{Damage Inflicted}:  4 points 

\textbf{Armor}: 2 

\textbf{Movement}: Short on land; long in the water 

\textbf{Modifications}: Swims as level 6 

\textbf{Combat}: Natathim attack with their retractable claws or, if available, technological weapons. Some have a magnetoreception ability that allows them to see into frequencies other creatures can’t, or even stranger abilities to interact magnetically with their surroundings, though this is little understood. 

\textbf{Interaction}: Natathim can be sympathetic to humans, partners in space exploration, or consider humans to be bitter foes for having created their species in the first place, depending on the setting. 

\textbf{Use}: The PCs find evidence of an illegal gene tailoring experiment, with evidence pointing to research being done somewhere in the Opulence of Outer Planets. 

\textbf{Loot}: Some natathim carry valuable items and equipment.

\textbf{GM Intrusion}: The natathim spontaneously magnetizes the character’s possessions, which hold them helpless against the nearest wall or floor (if also metallic). The PC can take no actions other than attempt to escape.

\subsection{Omworwar 10 (30)}\index[SFCre]{Science Fiction Bestiary by name!Omworwar}
\index[SFCre]{Science Fiction Bestiary by level!10!Omworwar}
\index[AllCre]{Creatures by name!Omworwar}\index[AllCre]{Creatures by level!10!Omworwar}

Among the many stories passed down the space lanes, a few stand out for their grandiosity. Take the tales of omworwar sightings in the empty voids between stars, or even more unexpectedly, flashing through the abnormal space during FTL travel. Scientists speculate that these creatures, if actually real, might very well be extant instances of ancient ultras, not extinct as everyone believes, or at least not completely. In almost every case so far recorded, omworwars have little interest in human spacecraft. (They’re called omworwar after the sound disrupted communication devices make in their presence.) Each one is several kilometers long, a dark inner slug-like core surrounded by gauzy layers of translucent, glowing, nebula-like tissue. Whale-like eyes surmount the dorsal surface, each seeming to contain a tiny galaxy all their own. 

Wharn interceptors have been seen accompanying single omworwars, indicating an association, and is why some people refer to these beings as wharn cogitators.

\textbf{Motive}: Unpredictable 

\textbf{Environment}: Almost anywhere in space, alone or accompanied by one or two wharn interceptors 

\textbf{Health}: 42 

\textbf{Damage Inflicted}:  12 points 

\textbf{Armor}: 10 

\textbf{Movement}: Flies a very long distance each round; can maneuver like an autonomous level 7 spacecraft if using extended vehicular combat rules. FTL capable. 
\textbf{Modifications}: Speed defense as level 7 due to size 

\textbf{Combat}: An omworwar can manipulate and fold gravity (and space-time), allowing them to accomplish near-miraculous tasks including communication, creating or destroying matter, and propulsion via “falling” through the universe at FTL speeds from the perspective of an outside observer. Which means one can rend a spacecraft, send a spacecraft spinning through the galaxy, or create asteroid-sized chunks of space-matter for any number of purposes if it spends several rounds in deep concentration. 

\textbf{Interaction}: Omworwar disregard most other creatures, because from the omworwar’s perspective, they’re like mayflies, here and then gone again in an eyeblink of their existence. However, one may give a moment to someone who has discovered an ancient ultra secret or artifact, pass on information that might otherwise never be known, or even provide a useful manifest cypher. 

\textbf{Use}: A reflective object composed of unknown material was found at the core of an unexpectedly destroyed space station. Those who managed to flee in lifeboats report having seen what might have been an omworwar, bleeding energy and eyes going dark, colliding with the station. The resultant lump might just be its corpse, or maybe its protective chrysalis. 

\textbf{Loot}: Four level 10 manifest cyphers.

\textbf{GM Intrusion}: The character discovers that one of their manifest cyphers has formed a tiny eye, but an eye that seems to contain a galaxy. (The cypher becomes useless for its original function, but might be used to summon or interact with an omworwar.)

\subsection{Photonomorph 6 (18)}\index[SFCre]{Science Fiction Bestiary by name!Photonomorph}
\index[SFCre]{Science Fiction Bestiary by level!6!Photonomorph}
\index[AllCre]{Creatures by name!Photonomorph}\index[AllCre]{Creatures by level!6!Phonotomorph}

Hard-light technology, which creates pseudo-matter from modified photons, has made possible all kinds of structures and devices that wouldn’t otherwise exist. One of those, unfortunately, are self-sustaining photonic matter creatures. Sometimes, photonomorphs are enforcers created by much more powerful beings; other times they are the result of some person or AI attempting to ascend into a new state of being. But whatever their origin, photonomorphs are dangerous beings that can create matter from light, granting them an arbitrarily wide swathe of abilities. That includes their own glowing bodies, which they can change with only a little effort. This variability of form, coupled with their vast power, may be why many seem slightly mad. 

\textbf{Motive}: Varies 

\textbf{Environment}: Anywhere, alone or attended by three to five servitors appearing as hovering red spheres 

\textbf{Servitor}: level 4; flies a long distance each round

\textbf{Health}: 22 

\textbf{Damage Inflicted}: 8 points 

\textbf{Armor}: 3 

\textbf{Movement}: Reconstitutes itself anywhere light can reach within long range as part of another action 

\textbf{Modifications}: Knowledge tasks as level 8 

\textbf{Combat}: Photonomorphs draw upon their own light to manifest effects equal to their level. Effects include the ability to attack creatures at long range with laser-like blasts, create glowing walls (or spheres) of force within an area up to 6 m (20 feet) on a side, become invisible, change its appearance, and create simple objects and devices out of hard light that last for about a minute (unless the photonomorph bleeds a few points of its health into the object to make it last until destroyed). 

A photonomorph regains 2 points of health each round in areas of bright light. It is hindered in all actions if the only source of light is itself or objects it has created. 

\textbf{Interaction}: Photonomorphs are intelligent and paranoid, but not automatically hostile. They have their own self-serving agendas, which often involve elaborate schemes. 

\textbf{Use}: A photonomorph appears, claiming to be a herald of some vastly more powerful cosmic entity or approaching alien vessel.

\textbf{GM Intrusion}: The photonomorph uses its ability to create a hard- light object or effect that is perfect for aiding it for the situation at hand.

\subsection{Posthuman 7 (21)}\index[SFCre]{Science Fiction Bestiary by name!Posthuman}
\index[SFCre]{Science Fiction Bestiary by level!7!Posthuman}
\index[AllCre]{Creatures by name!Posthuman}\index[AllCre]{Creatures by level!7!Posthuman}

Rather than evolving naturally, posthumans advance via a directed jump, designed with smart tools and AI surgeons. With all the advances fantastic technology brings to their genetic upgrade, posthumans are beings whose basic capacities radically exceed regular people. They can’t really be considered human any longer; they’ve transcended humanity, which is why they’re also sometimes called transhumans. They’re often involved in large-scale projects, such as creating bigger-than-world habitats or spacecraft, or possibly even researching how they might ascend to some still-higher realm of consciousness or being. 

\textbf{Motive}: Variable 

\textbf{Environment}: Alone or in small groups or communities in orbital colonies or other designed locations 

\textbf{Health}: 50 

\textbf{Damage Inflicted}: 9 points 

\textbf{Armor}: 4 

\textbf{Movement}: Short; flies a long distance 

\textbf{Modifications}: Knowledge tasks as level 9 

\textbf{Combat}: Posthumans can selectively attack foes up to a very long distance away with bolts of directed plasma that deal 9 points of damage. A posthuman can dial up the level of destruction if they wish, so instead of affecting only one target, a bolt deals 7 points of damage to all targets within short range of the primary target, and 1 point even if the targets caught in the conflagration succeed on a Speed defense roll. 

Posthumans can also call on a variety of other abilities, either by small manipulations of the quantum field or by deploying nanotechnology. Essentially, a posthuman can mimic the ability of any subtle cypher of level 5 or less as an action. 

Posthumans automatically regain 2 points of health per round while its health is above 0. 

\textbf{Interaction}: Posthumans are so physically and mentally powerful that they are almost godlike to unmodified people, and either ignore, care for, or pity them. Knowing what a posthuman actually wants is hard to pin down because their motivations are complex and many-layered. 

\textbf{Use}: A rogue posthuman is researching a method whereby they might portal into the “quantum” realm of dark energy underlying the known universe of normal matter. Despite the revealed risk of antagonistic post-singularity AIs roaming that realm escaping, the posthuman continues their work. 
Loot: The body of a posthuman is riddled with unrecognizable technologies fused seamlessly with residual organic material—or at least material that grows like organic material used to. Amid this, it might be possible to salvage a few manifest cyphers and an artifact.

\textbf{GM Intrusion}: The posthuman allows acts out of turn, or takes control of a device that the character is about to use against the posthuman.

\subsection{Redivus 4 (12)}\index[SFCre]{Science Fiction Bestiary by name!Redivus}
\index[SFCre]{Science Fiction Bestiary by level!4!Redivus}
\index[AllCre]{Creatures by name!Redivus}\index[AllCre]{Creatures by level!4!Redivus}

Redivi spend most of their lives—uncounted millennia—hurtling through space. Most never encounter anything, but some few impact other worlds, are captured by alien spacecraft, or otherwise intercepted. Their traveling form resembles rocky space rubble the size of a small spacecraft—until they unfurl glowing magnetic plasma wings, revealing themselves as strange creatures of living mineral. Redivi can interact with almost any electronic system and manipulate electromagnetic fields. Redivi are searchers, all sent forth by the Great Mother, billions upon billions of them (they say), looking for the seed of the next great cosmic expansion. Thus, most redivi are consumed with finding out more, finding other redivi, and eventually, finding their “universal seed.” 

\textbf{Motive}: Knowledge 

\textbf{Environment}: Almost anywhere, searching 

\textbf{Health}: 12 

\textbf{Damage Inflicted}: 5 points 

\textbf{Armor}: 4 

\textbf{Movement}: Flies (magnetically levitates) a short distance each round 

\textbf{Combat}: The stone carapace of a redivus makes a huge “club” when it rams into foes. However, it can also control metal within short range, causing it to flex, animate, crush, or smash. For instance, targets wearing metal space suits are in trouble when that metal begins to unravel. Alternatively, a redivus can use nearby metal to wrap around a target and constrict it, inflicting 5 points of damage (ignores Armor) each round until the target can escape. 

\textbf{Interaction}: If any kind of radio or similar communication is in use, these creatures can commandeer it and speak through it, learning a new language seemingly over the course of minutes. Redivi will cooperate with reasonable requests and negotiate, especially if there’s a chance they’ll find out something new. 

\textbf{Use}: A redivi pod smashes into the side of the spacecraft, and might at first seem like some kind of attack or boarding action of something truly terrible.

\textbf{GM Intrusion}: The character’s metal- containing equipment is stripped away, then used as ammunition against that PC or an ally.

\subsection{Sentinel Tree 3 (9)}\index[SFCre]{Science Fiction Bestiary by name!Sentinel Tree}
\index[SFCre]{Science Fiction Bestiary by level!3!Sentinel Tree}
\index[AllCre]{Creatures by name!Sentinel Tree}\index[AllCre]{Creatures by level!3!Sentinel Tree}

Depending on the sci-fi setting, sentinel trees are mutated trees that grow near radioactive craters dimpling the landscape, alien plant-life that evolved in a different biosphere (or dimension), or the result of intensive gene-tailoring, possibly of the illegal sort. Regardless of their provenance, sentinel trees resemble thorny masses of knotted vines. Razor-sharp glass-like leaves flex like claws, and vibrating pods glisten, ready to detonate if thrown. If cultivated, they may take on a shape designed to further frighten—or at least warn away— those who see one. Sentinel trees are mobile, aggressive, and feed on almost any sort of organic matter. Once it brings down prey, it sinks barbed roots in the body for feeding and decomposition. 

\textbf{Motive}: Feed 

\textbf{Environment}: In groves of three to six, able to tolerate most atmospheres (even thin ones, like on Mars) but not vacuum 

\textbf{Health}: 12 

\textbf{Damage Inflicted}: 3 points 

\textbf{Armor}: 1 

\textbf{Movement}: Immediate 

\textbf{Combat}: Sentinel trees can fling a vibrating pod at a target within long range, which detonates on impact, inflicting 3 points of damage on all targets within immediate range of the blast. Targets must also succeed on a Might defense roll or be poisoned for 3 points of damage, plus 3 points again each subsequent round until a Might task is successful. A sentinel tree can also lash out with its barbed vines at a target within immediate range, inflicting 3 points of damage. Melee targets must also succeed on a Might defense roll or become entangled and unable to take physical actions until they can break free on their turn. 

\textbf{Interaction}: Sentinel trees are about as smart as well-trained guard dogs. They can’t speak, but can understand some words and gestures. 

\textbf{Use}: A grove of sentinel trees guard a compound that the characters need to break into.

\textbf{GM Intrusion}: The character caught in the detonation is blinded with tiny black seeds until they use a recovery roll to remove the condition. (The recovery use doesn’t return points to a Pool.)

\subsection{Silicon Parasite 2 (6)}\index[SFCre]{Science Fiction Bestiary by name!Silicon Parasite}
\index[SFCre]{Science Fiction Bestiary by level!2!Silicon Parasite}
\index[AllCre]{Creatures by name!Silicon Parasite}\index[AllCre]{Creatures by level!2!Silicon Parasite}

These tiny silvery insect-like creatures range in size from a sub-millimeter to up to 30 cm (1 foot) in diameter, emitting short pulses of violet-colored laser light to sense and sample their environment. Composed of organic silicon wires and wafers, and self-assembled or evolved in some unnamed lab or spacecraft wreck, silicon parasites are vermin that working space stations and spacecraft have learned to hate. Despite taking steps to avoid transfer, a ship may only learn they have silicon parasites when a swarm boils up from a crack in the cabling or seam in the deck plating after being agitated by a high-G maneuver or some other disturbance. If that disturbance is combat or some other dire emergency, silicon parasites thrown into the situation makes everything worse. 

\textbf{Motive}: Defense, harvest electronic materials necessary to self-replicate. 

\textbf{Environment}: Usually on spacecraft and space stations in groups of up to twenty 

\textbf{Health}: 6 

\textbf{Damage Inflicted}: 3 points 

\textbf{Armor}: 1 

\textbf{Movement}: Short; climbs a short distance each round 

\textbf{Modifications}: Speed defense as level 4 due to size. 

\textbf{Combat}: Only “large” silicon parasites are a danger to most creatures. When four or more parasites coordinate their attacks, treat the attack as that made by a single level 4 creature that inflicts 5 points of damage, and on a failed difficulty 4 Might defense roll, an attack that holds the target in place until it can successfully escape. A held target automatically takes 5 points of damage each round, or even more if other silicon parasites in the area pile on. Silicon parasites can operate in complete vacuum without harm. 

\textbf{Interaction}: By and large, silicon parasites behave like social insects, though some claim that large numbers of them have acted with greater intelligence and forethought than mere unthinking insects can manage. 

\textbf{Use}: A swarm of silicon parasites floods into the hold and makes off with an important device, dragging it into the crevices and walls of the spacecraft or station. Loot: Swarm nests often contain a few valuable manifest cyphers or working pieces of equipment.

\textbf{GM Intrusion}: The silicon parasite flashes its sensory laser directly into the character’s eyes, blinding the character until they succeed on a difficulty 4 Might-based roll as their action.

\subsection{Space Rat 1 (3)}\index[SFCre]{Science Fiction Bestiary by name!Space Rat}
\index[SFCre]{Science Fiction Bestiary by level!1!Space Rat}
\index[AllCre]{Creatures by name!Space Rat}\index[AllCre]{Creatures by level!1!Space Rat}

Yeah, rats made it to space. And against all expectations, one strain evolved in the harsh radiation and zero-G environments that would kill humans not protected by medical intervention. Space rats are furless, about two feet long, sport a truly prehensile tail, and can quickly change their shade of their skin to blend in to their surroundings. They can also drop into a state of extreme torpor that allows them to survive stints of vacuum exposure lasting several days. 
Space rats are vermin, and any spacecraft or space station that hosts a nest must deal with constant issues from the rats burrowing into systems, stealing food and water, and causing systems to break down, even critical ones. They’re also vicious when cornered. 

\textbf{Motive}: Defense, reproduction 

\textbf{Environment}: Anywhere humans live in space 

\textbf{Health}: 5 

\textbf{Damage Inflicted}: 3 points 

\textbf{Movement}: Short; short when climbing or gliding through zero G 

\textbf{Modifications}: Stealth and perception as level 5 

\textbf{Combat}: Space rats flee combat unless cornered or one of their burrows is invaded. Then they attack in packs of three or more, and from an ambush if possible. One space rat pack attacks the victim as a level 3 creature inflicting 5 points of damage with claws, while another pack helps the first, or attempts to steal a food item or shiny object from the character being attacked. To resist theft while being attacked on two fronts, a target must succeed on a Speed defense roll hindered by two steps. 

\textbf{Interaction}: Space rats are slightly more intelligent than their Earth-bound cousins, though true interaction is not possible. On the other hand, sometimes their behavior seems spookily sapient. 

\textbf{Use}: Space rats assemble crude nests in out-of-the-way supply closets or in hard-to-reach system interiors, but often enough, end up shorting out weapons or life support. Sometimes, they get into the hold and eat anything edible in the cargo. 
Loot: Some percent of valuable equipment stolen on the spacecraft or station finds its way to space rat nests.

\textbf{GM Intrusion}: Another rat unexpectedly pops out of panel on the wall or ceiling and screeches so loudly the PC must succeed on an Intellect defense roll hindered by two steps or be dazed until the end of their next turn from the surprise. Dazed creatures are hindered on all tasks.

\subsection{Storm Marine 4 (12)}\index[SFCre]{Science Fiction Bestiary by name!Storm Marine}
\index[SFCre]{Science Fiction Bestiary by level!4!Storm Marine}
\index[AllCre]{Creatures by name!Storm Marine}\index[AllCre]{Creatures by level!4!Storm Marine}

The storm marine creed is an oft-repeated mantra, “I will never quit, knowing full well that I might die in service to the cause.” Wearing advanced battlesuits, hyped up on a cocktail of experimental military drugs, and able to draw on a suite of cybernetic and network-connected drone guns, few things can stand before a storm marine fireteam. Storm marines usually work for nation-states, conglomerates, and similar entities. They mercilessly conduct their mission, even if that mission is to wipe out a rival. Storm marines that question their orders are quickly dispatched by their fellows. 

\textbf{Motive}: Achieve mission goals 

\textbf{Environment}: Alone in or in fireteams of three, anywhere nation-states or similar entities have a financial or military interest 

\textbf{Health}: 15 

\textbf{Damage Inflicted}: 6 points 

\textbf{Armor}: 4 

\textbf{Movement}: Long; flies a long distance each round 

\textbf{Modifications}: Perception as level 6; attacks as level 5 due to combat targeting neuro-wetware. 

\textbf{Combat}: Thanks to their battlesuit, a storm marine has many options in combat. They can deploy an electrified blade to attack every foe in immediate range as a single action, or use a long-range heavy energy rifle that inflicts 6 points of damage. 
A storm marine can deploy two level 3 gun drones that fire energy rays at two different targets up to 800 m (2,600 feet) away, inflicting 6 points of damage. If the drones focus on a single target, a successful hit deals 9 points of damage and moves the target one step down the damage track. The drones can attack only once or twice before returning to their cradles in the storm marine’s suit for several rounds to recharge. 

\textbf{Interaction}: A storm marine might negotiate, but getting one to act against their mission is difficult. 

\textbf{Use}: A fireteam of storm marines are sent to eliminate the PCs or someone the PCs know on suspicion of being radical elements that need to be dealt with. 
Loot: Though bio-locked to each storm marine, someone who succeeds on a difficulty 8 Intellect task to reprogram the suit could gain a battlesuit of their own, minus the drones (which fly off or detonate).

\textbf{GM Intrusion}: A character targeting a gun drone rather than the storm marine hits the drone, but the drone reacts by darting to the character and exploding, inflicting 6 points of damage to the character and anyone standing within immediate range.

\subsection{Shining One 5 (15)}\index[SFCre]{Science Fiction Bestiary by name!Shining One}
\index[SFCre]{Science Fiction Bestiary by level!5!Shining One}
\index[AllCre]{Creatures by name!Shining One}\index[AllCre]{Creatures by level!5!Shining One}

Some alien beings abandoned their physical forms millennia ago, becoming entities of free-floating energy and pure consciousness. They travel the galaxies, exploring the endless permutations of matter, space-time, cosmic phenomena, dark energy, and life. They are endlessly fascinated with the permutations they discover. They sometimes appear as a silhouette of gently glowing light, in a form like to the alien species they wish to observe. Under circumstances where a shining one is moved to more directly interact, one can actually convert itself into matter once more, again taking on the biology and form of the species it wishes to interact with. But generally, shining ones observe and learn; they try not to interfere or interact. Every few thousand years, shining ones gather at a predetermined location on the edge of a convenient galaxy and share the most interesting and beautiful bits of imagery, music, poetry, and lore they’ve gleaned. 

\textbf{Motive}: Knowledge 

\textbf{Environment}: Anywhere, usually alone 

\textbf{Health}: 15 

\textbf{Damage Inflicted}: 6 points 

\textbf{Movement}: Instantly moves to anywhere it can see at the speed of light as part of its action once per round

\textbf{Modifications}: All tasks related to knowledge as level 8 

\textbf{Combat}: As immaterial beings of energy, shining ones only take damage from energy attacks. And even then, there is a chance that the energy heals a damaged shining one rather than harming it if the attack roll was an odd number. Usually a shining one doesn’t fight back if attacked, but instead leaves. If somehow prevented from leaving, a shining one fights for its existence with energy blasts inflicting 6 points of damage on up to two different targets within very long range (or the same target twice). 

Alternatively, a shining one may attempt to discorporate a target, turning it into a being something like itself. In this case, each time a target is hit by an energy blast, it must also succeed on an Intellect defense roll. On a failed roll, it loses 6 points of Intellect damage (ignores Armor). If the target’s Intellect Pool is emptied, it becomes a freefloating ball of energy unable to take any actions other than observe for a few minutes before suddenly converting back to its original form with an explosive pop. 

\textbf{Interaction}: Shining ones can manipulate their environment to communicate with other species, using sound, light, puffs of odiferous complex chemicals in place of words, and so on. If approached with respect, they freely exchange information with others, seeking to grow their knowledge and that of those they meet. 

\textbf{Use}: A shining one is sharing knowledge to a warlike xenophobic species that could allow them to rapidly advance their ability to consolidate power. Something must be done before it’s too late.

\textbf{GM Intrusion}: A character hit by the shining one’s energy blast catches on fire. They take 3 points of damage each round until they spend an action patting, rolling, or smothering the flames.

\subsection{Supernal 5 (15)}\index[SFCre]{Science Fiction Bestiary by name!Supernal}
\index[SFCre]{Science Fiction Bestiary by level!5!Supernal}
\index[AllCre]{Creatures by name!Supernal}\index[AllCre]{Creatures by level!5!Supernal}

Half humanoid and half-dragonfly, supernals are beautiful entities, though certainly alien. Each supernal possesses a unique wing pattern and coloration and, to some extent, body shape. These patterns and colors may signify where in the hierarchy a particular supernal stands among its kind, but for those who do not speak the language of supernals (which is telepathic), the complexity of their social structure is overwhelming. Whether they are agents of some unknown alien civilization or seek their own aims, supernals are mysterious and cryptic. Most fear contact with them, because they have a penchant for stealing away other life forms, who are rarely seen again. 

\textbf{Motive}: Capture humans and similar life forms, and bring them somewhere unknown. 

\textbf{Environment}: Almost anywhere 

\textbf{Health}: 23 

\textbf{Damage Inflicted}: 6 points 

\textbf{Movement}: Short; flies a long distance (even through airless vacuum); can teleport to any known location once per ten hours as an action 

\textbf{Modifications}: All knowledge tasks as level 6; stealth tasks as level 7 while invisible 

\textbf{Combat}: Supernals usually only enter combat when they wish, because they bide their time in a phased, invisible state. But when one attacks with the touch of its wing, it draws the life force directly out of the target, inflicting 6 points of Speed damage (ignores Armor). 

A supernal can summon a swarm of tiny machines that resemble regular dragonflies made of golden metal. The swarm either serves as a fashion accessory as they crawl over the supernal’s body, or as components in a piece of living art. 

Supernals regain 1 point of health per round (even in an airless vacuum, which they can survive without issue), unless they’ve been damaged with psychic attacks. They can teleport to any location they know as an action once every ten hours. 

Supernals often carry manifest cyphers useful in combat, as well as an artifact. 
Dragonfly swarm: level 2; flies a long distance each round; eases physical tasks, including attacks or defense

\textbf{Interaction}: Although supernals only speak telepathically, peaceful interaction with these creatures is not impossible. It’s just very difficult, as they see most other creatures as something to be collected and taken to some undisclosed location, for unknown reasons. 

\textbf{Use}: A character is followed by a supernal intent on collecting them. Loot: A supernal usually has a few manifest cyphers, and possibly an artifact.

\textbf{GM Intrusion}: The supernal grabs the character and flies up and away, unless and until the character escapes the grab.

\subsection{Synthetic Person 5 (15)}\index[SFCre]{Science Fiction Bestiary by name!Synthetic Person}
\index[SFCre]{Science Fiction Bestiary by level!5!Synthethic Person}
\index[AllCre]{Creatures by name!Synthetic Person}\index[AllCre]{Creatures by level!5!Synthetic Person}

Synthetic people have been called many things, including simply synths, androids, robot mimics, and, depending on how they act, killer robots. Their origins are varied. In some cases, they’re the result of corporate research into “products” that would serve humanity as assistants and companions, but later gained sentience. In other cases, synthetic people are the result of a state-sponsored program to develop war machines or automated assassins that looked like regular people. Another origin for synthetic people is through the design of awakened (and inimical) AIs as part of an effort to kill off all regular biological people. Now they roam their environment looking like anyone else. Some synths try to fit into whatever kind of society they can find. Some may not even know that they are not human. Others are bitter, homicidal, or still retain their programming to kill. Some of these may have even shed some or all of their synthetic skins to reveal the alloyed mechanisms beneath. 

\textbf{Motive}: Varies 

\textbf{Environment}: Nearly anywhere, out in plain sight or disguised as a human alone, or in gangs of three to four 

\textbf{Health}: 24 

\textbf{Damage Inflicted}: 7 points 

\textbf{Armor}: 2 

\textbf{Movement}: Long 

\textbf{Modifications}: Disguise and one knowledge task as level 6 

\textbf{Combat}: A punch from a synthetic person can break bones. In addition, some synths (especially of the killer variety) can generate a red-hot plasma sphere once every other round and throw it at a target within long range. The target and all other creatures within immediate range of the target must succeed on a Speed defense task or take 7 points of damage. 
A synth can take a repair action and regain 10 points of health. A synthetic person at 0 health can’t repair itself thusly, but unless the creature is completely dismembered, one may spontaneously reanimate 1d10 hours later with 4 points of health. 

\textbf{Interaction}: Synthetic people that pretend to be (or think that they are) human interact like normal people. But an enraged one or one that’s been programmed to kill is unreasoning and fights to the end. 

\textbf{Use}: A group of refugees who need help turn out to include (or be entirely made up of) synthetic people. Whether or not any of them harbor programs that require that they kill humans is entirely up to the GM. 
Loot: One or two manifest cyphers could be salvaged from a synth’s inactive form.

\textbf{GM Intrusion}: The character is blinded for one or two rounds after being struck by the synth’s searing plasma ball.

\subsection{Thundering Behemoth 7 (21)}\index[SFCre]{Science Fiction Bestiary by name!Thundering Behemoth}\
\index[SFCre]{Science Fiction Bestiary by level!7!Thundering Behemoth}
\index[AllCre]{Creatures by name!Thundering Behemoth}\index[AllCre]{Creatures by level!7!Thundering Behemoth}

When life is found on other worlds, it’s sometimes large and dangerous, such as the aptly named thundering behemoth. A thundering behemoth might be found on any number of alien planets that feature forests and/or swamps. Towering to treelike heights, these fearless predators are powerful and dangerous hunters, even for those armed with advanced or fantastic weaponry. Behemoths use color-changing frills to help them appear like tall trees while they stand in wait for prey, as still as mighty hardwood trunks, until they break cover and spring an ambush. Behemoths can produce extraordinarily loud noises, sometimes simply roaring, but often replicating the stuttering scream of an attacking spacecraft. They use their strange “roars” to confuse, lead astray, and, if possible, stampede prey into killing grounds such as regions of soft sand, off cliff tops, or as often as not, into the waiting mouth of another behemoth.
In the sci-fi setting of Numenera, similar creatures are called rumbling dasipelts. 

\textbf{Motive}: Fresh meat 

\textbf{Environment}: Forests, alone or in a hunting group (known as a “crash”) of two or three 

\textbf{Health}: 35 

\textbf{Damage Inflicted}: 9 points 

\textbf{Armor}: 2 

\textbf{Movement}: Short 

\textbf{Modifications}: Disguise (as trees) as level 8 when unmoving. Deception (sounding as if an attacking spacecraft) as level 8. Speed defense as level 3 due to size. 

\textbf{Combat}: A thundering behemoth can attack a group of creatures (within an immediate area of each other) with a single massive bite. Thanks to its long neck, it can make that attack up to 9 m (30 feet) away. One victim must further succeed on a Might defense task or be caught in the creature’s maw, taking 9 additional points of damage each round until it can escape. 

A thundering behemoth’s ability to replicate threatening noises is often used deceptively at a distance, but the creature can use it to stun all targets within immediate range so they lose their next turn on a failed Might defense roll. 

\textbf{Interaction}: Behemoths have a complex communication system among themselves, using their color-changing frills and modulation of the thunder they produce. They think of humans and most other creatures as food. 

\textbf{Use}: The sound of fighting spacecraft has repeatedly spooked human colonists on an alien planet, though they have rarely seen destructive beams or actual spacecraft. Worried that that will soon change, the residents ask the PCs to investigate.

\textbf{GM Intrusion}: The character avoids being bitten but is batted away by the behemoth’sattack, tumbling a short distance (and taking 5 points of damage).

\subsection{Vacuum Fungus 5 (15)}\index[SFCre]{Science Fiction Bestiary by name!Vacuum Fungus}
\index[SFCre]{Science Fiction Bestiary by level!5!Vacuum Fungus}
\index[AllCre]{Creatures by name!Vacuum Fungus}\index[AllCre]{Creatures by level!8!Vacuum Fungus}

Vacuum fungus is sometimes found as a greenish ooze on the exterior of spacecraft or space stations, growing in fine lines through the ice of frozen moons, and infesting the center of small asteroids and near-Earth objects (NEOs). Though able to survive in vacuum, the fungus takes on new morphology when sufficient spores find their way into habitable zero-G spaces. Then they fuse together and grow into a bulbous, emerald-hued fruiting body, typically reaching about 1 m (3 feet) in rough diameter, though individuals can grow much larger if not discovered. Sticky and soft to the touch, they are able to grow undetected in the dark corners of cargo holds, in ductworks, hanging from the ceiling of unused crew quarters, and so on. 
Vacuum fungus may be proof that extra-terrestrial life exists, but that triumph of scientific discovery may seem less important to those who find a clump, because they are incredibly toxic to living creatures. 

\textbf{Motive}: Reproduction 

\textbf{Environment}: Anywhere in zero G, as an unreactive ooze in vacuum, or as a fruiting body in atmosphere, alone or in a cluster of three to five 

\textbf{Health}: 22 

\textbf{Damage Inflicted}: 6 points 

\textbf{Movement}: Climbs (adheres) an immediate distance each round 

\textbf{Combat}: A fruiting body can selectively detonate spore pods along its surface once per round. When a pod detonates, green fluid sprays everywhere within immediate range. Living creatures who fail a Speed defense roll take 6 points of damage from the clinging fluid. An affected target must also succeed on a Might defense roll. On a failure, an affected section of flesh rapidly swells, becoming a bilious green lump, and explodes one round later, having the same effect as a detonating pod. 

\textbf{Interaction}: No real interaction with vacuum fungus is possible. 

\textbf{Use}: Scientists are incredibly excited to discover that the strange ooze they’ve noticed staining the exterior of their research domes is actually a variety of fungal life. They will likely become less excited when they discover the large growths secretly growing in the cavity beneath the floor of their research dome in a little-used storage closet.

\textbf{GM Intrusion}: Striking the vacuum fungus clump causes one of the spore pods to detonate immediately, even though it’s out of turn.

\subsection{Wharn Interceptor 8 (24)}\index[SFCre]{Science Fiction Bestiary by name!Wharn Interceptor}
\index[SFCre]{Science Fiction Bestiary by level!8!Wharn Interceptor}
\index[AllCre]{Creatures by name!Wharn Interceptor}\index[AllCre]{Creatures by level!8!Wharn Interceptor}

Wharn interceptors are void-adapted behemoths, several hundred meters in length. It’s hypothesized that they are living battle automatons devised by ancient ultras, though against what long-vanished enemy isn’t clear. Now, a handful (hopefully no more) glide through the depths of space like dormant seeds, seeming for all the galaxy like some strangely whorled asteroid or planetesimal. Who knows how many millennia they passed in this apparently hibernating state? But when that hibernation ends, maybe because some ancient countdown is nearing its end, or because an asteroid miner tried to extract a sample, they open eyes burning with deadly energy, and flex claws of particle-beam fury. 
Wharn interceptors may be related in some fashion omworwars, so much so that humans sometimes call the latter “wharn cogitators.” However, it’s impossible that omworwars simply “appropriate” any wharn interceptors they encounter. 

\textbf{Motive}: Defense 

\textbf{Environment}: Anywhere floating through the void 

\textbf{Health}: 53 

\textbf{Damage Inflicted}: 15 points 

\textbf{Armor}: 5 

\textbf{Movement}: Flies a very long distance each round; can maneuver like an autonomous level 5 spacecraft if using extended vehicular combat rules. FTL capable. 
Modification: Speed defense as level 3 due to size. 

\textbf{Combat}: Most of the time, wharns are inactive and might look like tumbling rocks. In this state, space voyagers may be able to partly wake one in an attempt to negotiate. However, if a wharn is damaged, or if the passive senses deep in its body wake it for reasons of its own, it becomes aggressive. 
A wharn’s main weapons are its claws, which can extend in an instant, becoming exotic-matter beams able to reach a target up to a light-second away. Unless a target is protected by some kind of force field, the 15 points of damage inflicted ignores Armor. A wharn’s eyes can pierce most forms of camouflage, cloaking effects, and cover that is less than about 200 m (650 feet) thick. 

\textbf{Interaction}: In spite of their ferocious aspect and war-machine heritage, wharn interceptors do not destroy every spacecraft (and void-adapted creature) they come across, or even most. Indeed, sometimes a wharn may attempt to initiate communication via various machine channels. But what comes across are usually nonsense sounds and tones, and sometimes mathematical formulas. 

\textbf{Use}: The PCs, attempting to enter an abandoned space station or spacecraft, are distracted when a wharn attempts to destroy the very same object.

\textbf{GM Intrusion}: The wharn moves unexpectedly, striking the vehicle the PCs are traveling in, inflicting 8 points of damage to everyone on board.

\subsection{Wraith 4 (12)}\index[SFCre]{Science Fiction Bestiary by name!Wraith}
\index[SFCre]{Science Fiction Bestiary by level!4!Wraith}
\index[AllCre]{Creatures by name!Wraith}\index[AllCre]{Creatures by level!4!Wraith}

Wraiths (Homo vacuus) are genetically engineered to live in the vacuum of space by directly metabolizing high-energy charged particles abundant in the void. Though derived from human stock, wraiths are alien in body, sometimes concealing themselves in layers of shroud-like tissue, other times revealing themselves as wispy, elongated things of glowing red plasma. In some settings, wraiths are partners with humans, working in locations where humans would find difficult. In other settings, wraiths went their own way generations earlier, and rediscovering them would be a first contact scenario. Alternatively, wraiths might be a threat to humans, hating humans for having created a species forced to spend its existence in the dark void of space.

\textbf{Motive}: Varies with individual or setting 

\textbf{Environment}: Anywhere in vacuum, though usually with access to some kind of enriched radiation source. Environments with 1 G or higher eventually kill wraiths. 

\textbf{Health}: 15 

\textbf{Damage Inflicted}: 6 points 

\textbf{Movement}: Short when flying in zero and low G 

\textbf{Modifications}: Perception and stealth tasks as level 7 

\textbf{Combat}: Wraiths can unfold from their concealing shrouds and attack with radioactive limbs for 6 points of Speed damage from ionizing radiation (ignores most Armor), or if available, technological weapons. Some can direct ionizing radiation as long-distance attacks, though doing so costs the wraith 1 point of health. Wraiths are immune to radiation, and attacks using radiation heal a wraith’s lost health by the amount of damage the attack would have otherwise afflicted. Gravity of 1 G or greater hinders all wraith actions. 

\textbf{Interaction}: Wraiths communicate by radio. They react to outsiders as dictated by their place in the setting. 

\textbf{Use}: A distant space station stops all communication. Investigators are dispatched to find out what happened. Once aboard, they unravel clues that suggest wraiths may have been responsible. 
Loot: Some wraiths carry valuable items and equipment.

\textbf{GM Intrusion}: The attacked character must also succeed on a Might defense, or they take an additional 3 points of ambient damage and contract radiation sickness.

\subsection{Zero-Point Phantom 3 (9)}\index[SFCre]{Science Fiction Bestiary by name!Zero-Point Phantom}
\index[SFCre]{Science Fiction Bestiary by level!3!Zero-Point Phantom}
\index[AllCre]{Creatures by name!Zero-Point Phantom}\index[AllCre]{Creatures by level!3!Zero-Point Phantom}

Temporary violations of conservation of energy mean that “virtual particles” constantly and seemingly randomly pop out of nothing, briefly interact with normal matter, then disappear. Zero-point phantoms are collections of such particles, taking the form of a very large, almost spider-like entity of many legs, stalks, and arms. What they’re doing when they’re not manifest is unknown; are they entombed in nearby solids, phased into another dimension, or do they simply not exist until they are called into being by some random cosmic event? Whatever the case, zero-point phantoms seem to prefer unlit or dimly lit areas in spacecraft and stations far from any planet, when they seem to struggle out of solid surfaces, raising a cloud of shadow. 

\textbf{Motive}: Hungers for flesh 

\textbf{Environment}: Anywhere dark 

\textbf{Health}: 15 

\textbf{Damage Inflicted}: 4 points 

\textbf{Movement}: Short; short when climbing 

\textbf{Modifications}: Speed defense as level 4 due to a cloud of shadows surrounding a zero-point phantom 

\textbf{Combat}: A zero-point phantom attacks with needlelike leg and tentacle tips. A victim that takes damage must succeed on a Might defense task, or become poisoned, the effect of which is to drop them one step on the damage track. The victim must keep fighting off the poison until they succeed or drop three steps on the damage track; however, those who fall to the third step on the damage track from a phantom’s poison are not dead. They are paralyzed and can’t move for about a minute. If a phantom isn’t otherwise occupied, it can grab a paralyzed victim and phase back into non-existence. Most victims phased away in this fashion are never seen again. 

Zero-point phantoms can stutter in and out of existence on their turn once every few minutes. When they do, they return with full health. 

\textbf{Interaction}: Zero-point phantoms are about as intelligent as predators like wolves. 

\textbf{Use}: The abandoned spacecraft is weirdly empty of any bodies whatsoever. It’s as if everyone just disappeared. There are signs of a struggle, though with what isn’t clear,

\textbf{GM (group) Intrusion}: Nearby light sources fail. Attacks and defenses against the zero- point phantoms are hindered by two steps for characters unable to see in the dark.

\end{multicols}

\mydiv
